\chapter*{\introductionname}
\label{cap:intro}
\addcontentsline{toc}{chapter}{\introductionname}

El aprendizaje automático ha avanzado significativamente a traves de los años. Áreas como la vision por computador, el procesamiento de lenguaje natural y la robótica han experimentado mejoras 
notables gracias a la aplicación de técnicas de aprendizaje profundo. Estas técnicas permiten a los modelos aprender representaciones complejas de datos, lo que ha llevado a avances en tareas 
como el reconocimiento de imágenes, la traducción automática y la conducción autónoma. Sin embargo las redes neuronales no son la bala de plata para todos los problemas. En particular, el modelamiento de fenómenos
físicos complejos la mecánica de fluidos sigue siendo un desafío debido a su naturaleza no lineal y la alta dimensionalidad de los sistemas involucrados. En este contexto, las redes neuronales informadas por la física (PINNs) han surgido 
como una herramienta prometedora para abordar estos desafíos, integrando principios físicos directamente en el proceso de aprendizaje.
El volcán Villarica, ubicado a 2847 m sobre el nivel del mar, es un estratovolcán parcialmente glaciado dentro de la Zona Volcánica del Sur (SVZ) de los Andes de Chile \cite{https://doi.org/10.1029/2018GC007692}.